% !TEX root = ../main.tex
\chapter{Architecture}\label{ch:architecture}
This chapter is for a researcher who wants to check what the app does, or an agent who wants to
build on it. It names the code for every claim. Paths are relative to the repository root.

\section{Two halves: the core and the web UI}\label{sec:architecture-halves}
The repository holds two programs that meet at one seam.

\begin{description}
\item[The core] \code{Working/} and \code{Adapters/}, plus the batch scripts in \code{Pipelines/}.
  Plain Python with numpy, scipy and friends. \code{Working/} holds configuration
  (\code{Working/config.py}), recipe building and hashing (\code{Working/recipes.py}), execution
  (\code{Working/execution.py}), the database layer (\code{Working/database/}), units
  (\code{Working/units.py}), registration, the Library, Discovery and training code, and the cluster
  job writers (\code{Working/hpc/job_export.py}, \code{Working/discovery/seed_job.py},
  \code{Working/training/}). \code{Adapters/} holds the block registry: one file per
  \term{block}{block}, each registering an \code{AdapterSpec} (\code{Adapters/base.py}).
\item[The web UI] \code{webui/}. A React and TypeScript client in \code{webui/client/} (built with
  Vite), and a FastAPI \emph{bridge} in \code{webui/server/}, started by \code{webui/run_server.py}.
  The client never touches the database or a file: every read and write is an HTTP call through
  \code{webui/client/src/api.ts} and the per-workspace modules in \code{webui/client/src/api/}. The
  bridge answers those calls by calling the core.
\end{description}

The dependency runs one way. The bridge imports the core; the core imports nothing from the bridge
and no UI library at all. This is rule~1 of \code{CLAUDE.md}: \code{Working/}, \code{Adapters/} and
\code{Pipelines/} import neither a browser library, nor FastAPI or uvicorn, nor the retired Panel
family (Panel, HoloViews, Bokeh). The one drawing library the core keeps is matplotlib, for figure
export.

\begin{plainterms}[Why the engine has no dashboard bolted to it]
Think of a car engine on a test bench. If the engine only runs when it is bolted into one particular
car, with that car's dashboard wired in, you cannot test it on the bench or put it in a different car.
The core is the engine. Because it never assumes a screen exists, the same code runs on the laptop
behind the web pages, on a computing cluster with no screen at all, and inside automated tests.
\end{plainterms}

Three things depend on that separation.
\begin{itemize}
\item \textbf{Cluster execution.} A long \term{chain}{chain} is written as a SLURM job that runs
  the same core on the university cluster (\cref{sec:architecture-long-work}). The cluster
  environment has no web server and no browser.
\item \textbf{Headless tests.} The \code{pytest} suite in \code{tests/} runs the core without
  starting the bridge or a browser.
\item \textbf{Reproducibility.} A result is a function of a recipe and the stored samples, not of
  what a page happened to show (\cref{sec:architecture-reproducibility}).
\end{itemize}

\begin{underthehood}[How the rule is enforced]
\code{tests/test_import_boundaries.py} parses every Python file in the repository with \code{ast}
and lists every module it imports, at module scope or inside a function. It fails if any file
imports \code{panel}, \code{holoviews}, \code{bokeh} or the retired \code{UI} package
(\code{test_nothing_in_the_repository_imports_the_panel_family},
\code{test_nothing_in_the_repository_imports_the_retired_ui_package}), and if any file under
\code{Working/}, \code{Adapters/} or \code{Pipelines/} imports \code{fastapi}, \code{uvicorn} or
\code{starlette} (\code{test_the_core_does_not_import_fastapi_or_uvicorn}). A fourth test checks
that the walk actually reached those folders, so the check cannot pass by walking nothing. The test
reads Python only; the TypeScript client is not walked.
\end{underthehood}

\section{What happens when you press Run}\label{sec:architecture-run}
This section follows one chain run in Analyse, from the button to the stored result.

\begin{enumerate}
\item \textbf{The client posts the chain.} \code{startRun} in \code{webui/client/src/api.ts} sends
  \code{POST /api/runs} with the recording id, the \term{span}{span} (or none, for the whole
  channel), the ordered steps (stage, algorithm, settings, \term{side inputs}{side-input}), the
  plot width in pixels, and whether to carry a \term{null}{null}.
\item \textbf{The bridge checks it.} \code{post_run} in \code{webui/server/app.py}:
  \begin{itemize}
  \item refuses an unknown recording (404) and the \term{held-out}{held-out} recording (423);
  \item checks every junction between steps with \code{chain.validate}
    (\code{webui/server/chain.py}): each block's input type must match what the step before
    produces (422 otherwise);
  \item builds the recipe with \code{chain.build_recipe}, which passes each step's settings through
    that block's own \code{validate_params} first (defaults filled in, types coerced, ranges checked),
    then calls \code{Working.recipes.make_recipe};
  \item refuses a span longer than any step's \code{max_span_samples} (422);
  \item estimates the run's time from each block's \code{estimate} and routes it
    (\code{chain.route_local}). A run over the Analyse local limit is not started: the bridge writes
    a SLURM job for it instead and answers 422 with the \code{sbatch} command
    (\code{_slurm_refusal}; \cref{sec:architecture-long-work}). A run whose cost model is still
    being calibrated on this machine is held with 409.
  \end{itemize}
\item \textbf{A job starts.} \code{JobManager.start} (\code{webui/server/jobs.py}, built on
  \code{RunManager} in \code{webui/server/runs.py}) starts one worker thread per run and writes a
  row to the \code{jobs} table. The thread calls \code{execute_recipe} unchanged, passing callbacks
  for progress, per-step results and cancellation. It always passes \code{force=True}, so the step
  loop runs and the step cache, not a reused run, supplies any unchanged prefix.
\item \textbf{The core runs the recipe.} \code{execute_recipe} in \code{Working/execution.py}
  opens the database through \code{init_db}, then, in \code{_execute_recipe_with_conn}:
  \begin{itemize}
  \item refuses the held-out recording unless \code{HELD_OUT_UNLOCK} is set in
    \code{Working/config.py};
  \item validates the steps again (\code{validate_recipe_steps}), because a recipe written by
    hand or on the cluster may not have come through \code{make_recipe};
  \item stores the recipe as a \code{configs} row (\code{get_or_create_config}) and inserts a
    \code{runs} row with status \code{running} before any step runs;
  \item memory-maps the channel's \code{.npy} file and slices out the span
    (\code{_load_signal}), so only the span is read from disk;
  \item for each step: checks for cancellation, records the current step on the run row, validates
    the settings, checks the span ceiling, then looks the step up in the step cache
    (\cref{sec:architecture-reproducibility}). On a hit it restores the typed output from disk; on a
    miss it calls the block's \code{run} with the signal, the time axis, the
    \term{sampling rate}{sampling-rate}, the settings, and whatever extra arguments the block's
    \code{run} declares (\code{value}, its side inputs, \code{conn} and \code{recording}, or
    \code{units}). A step slower than \code{STEP_CACHE_WRITE_THRESHOLD_S} (1\,s) is written to the
    cache;
  \item writes a \term{span set}{span-set} to the \code{detections} table, shifted from
    span-relative to channel-absolute sample indices. A block that takes a span set in and gives one
    out (a feature block) does not write it again;
  \item calls the block's \code{persist} hook, if it has one, which writes a file and returns its
    path; the path is recorded as an \code{artifacts} row;
  \item marks the run \code{completed} with its duration and per-step timings. Any exception marks
    it \code{failed} with the full traceback in \code{error_text}; a cancel marks it \code{failed}
    with text starting ``Cancelled''.
  \end{itemize}
\item \textbf{Results stream back.} Progress and each step's result go to the page as Server-Sent
  Events on \code{GET /api/runs/{job}/events}. Every event is also kept on the job, so a page
  reloaded mid-run replays it. Closing the page does not cancel the run. The client draws each
  step's output from \code{GET /api/runs/{job}/steps/{index}}, which serialises the typed value
  through one function, \code{to_payload} in \code{webui/server/serialize.py}.
\end{enumerate}

\begin{watchout}[Cancel is checked between steps]
\code{execute_recipe} checks for cancellation before each step. A block that declares
\code{should_cancel} on its \code{run} also checks inside its own loop; any other block runs its
step to the end before the cancel takes effect. This is why the bridge refuses, before it starts, a
run it estimates is over the local limit.
\end{watchout}

\section{Where things are kept}\label{sec:architecture-storage}
\subsection{The database}
The project's database is one SQLite file, \code{DATA/db/annotations.sqlite}
(\code{Working/database/schema.py}, \code{DB_PATH}). It is plain SQL through Python's
\code{sqlite3}; there is no ORM (rule~3 of \code{CLAUDE.md}).

The whole schema is in \code{Working/database/schema.py}. \code{init_db()} creates every table with
\code{CREATE TABLE IF NOT EXISTS} and adds new columns to existing tables with
\code{_migrate_columns}, which reads \code{PRAGMA table_info} and adds only what is missing. So
\code{init_db()} can run on every start, and a change to the schema is an addition. Two migrations
cannot be additive, because SQLite cannot alter a \code{CHECK} constraint in place:
\code{_migrate_annotations_verdict} and \code{_migrate_review_queues_kinds} rebuild their table. The
schema's docstring says the first backs up the file, runs in one transaction and checks row counts
before committing. One further migration rewrites rows: \code{_migrate_legacy_detections} shifts
detections written span-relative before 2026-09-21 to channel-absolute indices, and records each
rewrite in \code{audit_log}. \Cref{app:data} lists the tables a researcher meets.

\subsection{Bulk arrays live on disk}
Rule~4 of \code{CLAUDE.md}: a large array never enters the database. A recording's samples are a
\code{.npy} file per channel under \code{DATA/derived/channels/}; a matrix profile, a window
matrix, a cached step output, a model file are files too. The database row holds the path. For
example \code{recordings.npy_path} names the channel file, \code{artifacts.path} a run's output
file, and \code{step_artifacts.path} a cached step's directory.

\subsection{Machine rows and human rows}
Rule~5 of \code{CLAUDE.md}: \term{detections}{detection} are written only by the machine and
\term{annotations}{annotation} only by people, in separate tables, and no code path writes one into
the other. A person's judgement of a machine detection is an
\term{adjudication}{adjudication}, in its own table that points at the detection; the detection row
itself is never edited. \Cref{sec:core-ideas-detections} explains the idea in plain words.

\begin{underthehood}[The two doors]
\code{webui/server/writes.py} holds \code{write_human(conn, table, row)} and
\code{write_machine(conn, table, row)}. Each accepts only its own list of tables
(\code{HUMAN_TABLES}: annotations, adjudications, tags, templates, groupings, hand edits and the
\code{motif_*} Library tables; \code{MACHINE_TABLES}: runs, configs, detections, artifacts,
encodings, the step cache, recordings, registered artifacts, window sets, Discovery's sessions and
runs) and refuses any other with a \code{PermissionError} naming the rule. \code{settings} and
\code{audit_log} are on neither list. The executor itself writes \code{detections} through
\code{Working.database.runs.insert_detection}, not through these doors.
\end{underthehood}

\section{Reproducibility: the same inputs give the same run}\label{sec:architecture-reproducibility}
\subsection{The recipe and its hash}
Every run is described by a \term{recipe}{recipe} (\code{Working/recipes.py}): the recording id,
the span (or none for the whole channel), and the ordered steps, each with its stage, block name,
settings and side-input bindings, plus an optional fan-out over channels or bands.
\code{canonical_json} writes it with sorted keys, no whitespace variation, and numpy numbers
turned into plain Python numbers; \code{recipe_hash} is the SHA-256 of that text, and
\code{short_hash} its first eight hexadecimal characters. Two recipes that mean the same thing get
the same hash whatever order their keys were written in.

Because the bridge fills each block's defaults (\code{validate_params}) before hashing, a chain
whose settings were left at their defaults and the same chain with the defaults typed in get one
hash. A side input bound to a Library exemplar is hashed by what it is (source file, channel, sample
range), not by its row id in this database (\code{_normalize}).

\begin{watchout}[The recording is named by its row id]
The recipe names its recording by \code{recording_id}, the row id in this database's
\code{recordings} table. The same recipe sent to another copy of the database, where the same
channel has a different id, hashes the same but means a different channel. The cluster job
writers carry the database file or the channel paths with the job for this reason
(\cref{sec:architecture-long-work}).
\end{watchout}

\subsection{Run reuse and the step cache}
The short hash is the key of the \code{configs} row (\code{get_or_create_config}). A run row links
a config to a recording and a span, so the provenance of any result is the run row, its config (the
full recipe as JSON), its timings and error text, its detections and its artifacts.

\code{execute_recipe} uses the hash twice.
\begin{description}
\item[Run reuse] Called without \code{force}, it returns an existing completed run of the same
  config over the same recording and span instead of recomputing (\code{find_completed_run}). The
  bridge always passes \code{force=True}, so this applies to scripts, not to the pages.
\item[The step cache] Each step is keyed by the hash of the recipe cut off after that step
  (\code{_recipe_prefix_hash}). Changing step $N$ changes the key of step $N$ and every later step,
  but no earlier one (\code{invalidated_step_indices}), so re-running a chain after editing a late
  step restores the early steps from disk. Entries are rows of \code{step_artifacts} pointing to a
  directory under \code{STEP_CACHE_ROOT} (\code{DATA/derived/step_cache} in
  \code{Working/config.py}); the typed output is written and read by its own type's
  \code{to_path}/\code{from_path}. Nothing is evicted.
\end{description}

Two exceptions keep the cache honest. A block whose \code{run} declares \code{conn} reads the live
database (for example, human verdicts); its output depends on more than the recipe, so that step
and every later one skip the cache, and run reuse is off for the whole recipe
(\code{_first_store_reading_step}). And where a step that is told the recording's unit sits in the
prefix, a unit other than volts enters the key, so a result computed under one unit is never served
for another (\code{_recipe_prefix_hash}).

\section{Units and missing values}\label{sec:architecture-units}
\subsection{Units}
The recordings are not all stored in one unit: some are volts, some millivolts. The unit is a
property of the data, in \code{recordings.units} (\code{V}, \code{mV} or \code{uV}; empty means
undeclared) with \code{recordings.units_note} saying where it came from. \code{Working/units.py} is
the vocabulary: \code{parse_units} reads a manifest's text, \code{to_mv_factor} gives the factor to
millivolts (\code{None} when undeclared, never a guessed 1), and \code{read_as_mv} tells a block how
to read its samples as millivolts.

\begin{plainterms}[Why the unit matters]
A ruler marked only with numbers is useless if you do not know whether they are centimetres or
inches. A threshold of 0.005 means five millivolts on a recording stored in volts and five
thousandths of a millivolt on one stored in millivolts: a factor of a thousand.
\end{plainterms}

The core does not convert samples. Every block receives them as stored. Two places use the unit:
\begin{itemize}
\item the bridge, when it draws a trace (\code{display_channel} in \code{webui/server/corpus.py});
\item a block that prints millivolts. Its \code{run} declares a \code{units} argument, and
  \code{execute_recipe}, the block page's read-out and the window chain pass it
  \code{recordings.units}. The block converts with \code{read_as_mv}, which reads an undeclared unit
  as volts and returns a note saying so, for the page to print beside the number. At the time of
  writing these are the drop detector's two blocks and \code{interrogation.event_shape}.
\end{itemize}
Any other block works in the stored unit, and any threshold in its settings means different things
on a volts recording and a millivolts one. Each block's entry in \cref{app:blocks} says which kind it
is. \code{init_db()} fills in the unit measured for each known file on 2026-09-23
(\code{RECORDING_UNITS_EVIDENCE}) on rows that carry no unit and no note, so a unit a person
declares later is never overwritten. A unit is declared in Settings \textrightarrow{} Datasets
(\apppage{settings.datasets}); \cref{app:data} has the details.

\subsection{Missing values}
A gap in a recording is stored as NaN (not a number). The executor does nothing special with NaN:
\code{_load_signal} slices the samples as they are and hands them on. What happens to a NaN is
therefore each block's decision, and each block's entry in \cref{app:blocks} traces it. Some parts
of the system refuse NaN outright: the Library's content hash raises an error on a NaN sample
(\code{Working.library.identity.content_hash}), and registering a matrix profile fails its
\code{finite} check (\code{docs/DATA_REGISTRATION.md}). For the drop detector the researcher chose to
skip gaps rather than refuse a recording that has them (2026-10-14);
\unverified{the same skip-the-gaps choice has not been made for every block}.

\section{The bridge's sandbox}\label{sec:architecture-sandbox}
The bridge runs in one of two modes, chosen when it starts and printed in its log
(\code{webui/server/runtime.py}, \code{Runtime}).

\begin{description}
\item[\code{--sandbox} (the default)] The database is copied into
  \code{webui/runtime/<stamp>/}. Before any run, every path the core could write is redirected
  under that folder by rebinding module attributes: \code{Working.config.STEP_CACHE_ROOT}, the
  results folders of the matrix-profile, window-matrix, clustering, trace-shape and shape-cluster
  blocks, the window-pool root, the classifier's model root, and the three cost-calibration files
  (matrix profile, block cost, window matrix).
  The step-cache write threshold is lowered to zero so every step is cached. The bridge refuses to
  start if any redirected path lies outside the runtime folder. Nothing you do in a sandbox reaches
  the real database. The smoke walk (\code{webui/smoke.py}) runs against a sandbox.
\item[\code{--project}] The real \code{DATA/db/annotations.sqlite} is opened in place, in WAL
  journal mode, after a timestamped backup is written to \code{DATA/db/backups/} (the last ten are
  kept). The caches and results folders stay where the core has them. This is the mode in which
  the researcher's annotations and runs accumulate.
\end{description}

In both modes recordings are only read, memory-mapped, and the held-out recording is refused on
every route. \code{webui/start.ps1} starts a sandbox unless given \code{-Project};
\code{webui/start.sh} takes \code{--project} as its first argument.

\begin{watchout}[Never point the app at a link to the real data]
\code{CLAUDE.md} records that removing a git worktree that held a directory junction into the
real \code{DATA/} folder deleted the real database and derived channels (2026-09-21). Remove such a
junction (\code{cmd /c rmdir}) before removing the worktree, and never run the bridge, the tests or a
block through one.
\end{watchout}

\section{Long work: jobs and the cluster}\label{sec:architecture-long-work}
\subsection{Jobs on this machine}
Any piece of work longer than a page should wait for is a \term{job}{job}
(\code{webui/server/jobs.py}). A chain run is a job of kind \code{chain_run}; other work starts
through \code{JobManager.start_job(kind, fn, meta)}, with kinds \code{sweep}, \code{import},
\code{regroup}, \code{training} and \code{cross_channel} (\code{KINDS}). Each job:
\begin{itemize}
\item runs in its own thread and is one row of the \code{jobs} table, updated as it progresses
  and when it ends;
\item streams progress as Server-Sent Events on \code{GET /api/jobs/{id}/events}, replayed in full
  to a late subscriber;
\item can be cancelled with \code{POST /api/jobs/{id}/cancel}, which sets a flag the job checks
  when it chooses (a chain run checks it between steps);
\item survives a restart of the bridge as a snapshot: a job no longer in memory is rebuilt from its
  row and, for a chain run, from the \code{runs} row the core wrote.
\end{itemize}

\subsection{Work sent to the cluster}
The site never logs in to the cluster or submits anything. It writes a job (a recipe or spec, an
\code{sbatch} script, and a list of what to copy), you copy the files and run \code{sbatch}, and you
bring the results back. \Cref{ch:jobs} describes this from the user's side; the code is:

\begin{description}
\item[Routing] \code{route_recipe} in \code{Working/hpc/job_export.py} sums each step's
  \code{estimate} over the span (times the fan-out width) and answers \code{local},
  \code{cluster} or \code{unknown} (a step's estimator is present but not calibrated on this
  machine). The default ceiling is \code{CLUSTER_ROUTING_CEILING_S}, 900\,s
  (\code{Working/config.py}). The Analyse limit can be set in Settings \textrightarrow{} Compute \&
  HPC (\code{limit.Analyse}, in minutes; \code{analyse_ceiling_s} in \code{webui/server/chain.py});
  Discovery reads its own \code{limit.Discovery} (\code{Working/discovery/fanout.py}).
\item[Writing the job] \code{export_job} writes the recipe as JSON and an \code{sbatch} script that
  runs it through \code{Pipelines/run_recipe/run_recipe.py}, which calls \code{execute_recipe}. The
  script's comments list the code modules, input files and outputs the job needs
  (\code{recipe_dependencies}). The wall time is the estimate times a safety factor, rounded up, and
  clamped to \code{HPC_MAX_WALLTIME_MINUTES} (20 minutes, this account's limit;
  \code{_slurm_time_from_estimate}). A job that needs longer than that is still written, and the
  bridge's message says it will most likely be killed. The script changes directory to
  \code{HPC_REMOTE_REPO_ROOT}; \code{Working/hpc/job_export.py}'s docstring warns this value should be
  confirmed for the account before submitting.
\item[Other writers] Models \textrightarrow{} Launch writes a CNN training job directory
  (\code{Working/training/shape_cnn.py}, \code{export_cnn_job}); the Shape clustering page writes a
  Ward clustering over every training window (\code{Working/training/full_ward.py}); Discovery's Seed
  page writes a seed search (\code{Working/discovery/seed_job.py}), which checkpoints its result file
  after every channel and resubmits itself while the file reads incomplete.
\item[Bringing results back] A returned job directory is imported by
  \code{Working.training.hpc_import.import_results}, the same function the command
  \code{python -m Working.training import-results} calls. It checks that the recipe hash in
  \code{recipe.json}, \code{job.json} and \code{out/done.json} agree, that the windows and pool match
  this database, and that every prediction is for a window of the job; a mismatch raises
  \code{ImportRefused} and nothing is stored. A seed search's result file is imported by Discovery's
  own import function. A matrix-profile job that runs through \code{run_recipe.py} registers its own
  rows through the block's \code{persist} hook. Recipe-and-script jobs written by Discovery's chain
  export and Models' paired training have no return path through the site.
\end{description}

\code{GET /api/hpc/cluster} (\code{webui/server/jobs_routes.py}) lists every job the site wrote,
whatever wrote it, with a state read from its folder, its result file and the database.
